% from ega0-20-fr.tex
\oldpage[0\textsubscript{IV-1}]{135}
$w_*:\operatorname{Exan}_B(C,\mathfrak{R}/\mathfrak{R}^2)
\to\operatorname{Exan}_B(C,L)$ correspondant à un $B$-homomorphisme
$w:\mathfrak{R}/\mathfrak{R}^2\to L$. Le fait que c) entraîne b) résulte
alors de la commutativité du diagramme (18.3.6.5)
\[
\xymatrix{
  \operatorname{Exan}_B(C,\mathfrak{R}/\mathfrak{R}^2)
    \ar[r]^{w_*} \ar[d]_{u^*} &
  \operatorname{Exan}_B(C,L) \ar[d]^{u^*} \\
  \operatorname{Exan}_A(C,\mathfrak{R}/\mathfrak{R}^2)
    \ar[r]_{w_*} & \operatorname{Exan}_A(C,L).
}
\]

\medskip
(iv) On a vu (20.1.7) que les inverses à droite de $E\to C$
correspondent biunivoquement de façon canonique à l'ensemble des éléments
$D\in\operatorname{D\acute{e}r}_A(E,\mathfrak{R}/\mathfrak{R}^2)$ tels que $D(x)=x$ dans
$\mathfrak{R}/\mathfrak{R}^2$, donc aussi, par (20.4.8), à l'ensemble des
$E$-homomorphismes
$h:\Omega^1_{E/A}\to\mathfrak{R}/\mathfrak{R}^2$ tels que le composé
\[
  \mathfrak{R}/\mathfrak{R}^2\xrightarrow{d_{E/A}}
  \Omega^1_{E/A}\xrightarrow{h}\mathfrak{R}/\mathfrak{R}^2
\]
soit l'identité. Par tensorisation avec $C$, on en déduit (puisque
$\mathfrak{R}/\mathfrak{R}^2$ est un $C$-module) que le composé
\[
  \mathfrak{R}/\mathfrak{R}^2\xrightarrow{\delta_{C/E/A}}
  \Omega^1_{E/A}\otimes_E C\xrightarrow{h\otimes1}
  \mathfrak{R}/\mathfrak{R}^2
\]
est l'identité; or, puisque $\mathfrak{R}/\mathfrak{R}^2$ est un
$C$-module, $h\mapsto h\otimes1$ est un isomorphisme de l'ensemble
$\operatorname{Hom}_E(\Omega^1_{E/A},\mathfrak{R}/\mathfrak{R}^2)$ sur
$\operatorname{Hom}_C(\Omega^1_{E/A}\otimes_E C,\mathfrak{R}/\mathfrak{R}^2)$; et par
ailleurs (ii) prouve que l'on peut identifier canoniquement
$\Omega^1_{E/A}\otimes_E C$ et $\Omega^1_{B/A}\otimes_B C$,
$\delta_{C/E/A}$ s'identifiant alors à $\delta_{C/B/A}$. C.Q.F.D.
\end{proof}

\begin{example}[20.5.13]
\label{0.20.5.13-fr}
Soient $B=A[X_\alpha]_{\alpha\in I}$ une algèbre de polynômes sur $A$,
$\mathfrak{R}$ un idéal de $B$, $(P_\lambda)$ un système de générateurs
de $\mathfrak{R}$ et $C=B/\mathfrak{R}$; on sait que
$\Omega^1_{B/A}$ est un $B$-module libre dont les $dX_\alpha$ forment une
base (20.4.13, (i)), donc les $dX_\alpha$ forment aussi une base du
$C$-module libre $\Omega^1_{B/A}\otimes_B C$. D'autre part, il résulte
aussitôt de la définition que l'image de
$\mathfrak{R}/\mathfrak{R}^2$ par $\delta_{C/B/A}$ est le sous-$C$-module
engendré par les
\[
  dP_\lambda=\sum_\alpha
  \frac{\partial P_\lambda}{\partial X_\alpha}\,dX_\alpha.
\]
On en conclut que $\Omega^1_{C/A}$ est isomorphe au quotient du
$C$-module libre ayant les $dX_\alpha$ pour base, par le sous-$C$-module
engendré par les $dP_\lambda$, ce qui donne une description d'un module
de différentielles d'une algèbre \emph{quelconque}, toute $A$-algèbre
$C$ pouvant s'obtenir de la façon précédente.
\end{example}

\begin{corollary}[20.5.14]
\label{0.20.5.14-fr}
Si $C$ est une $A$-algèbre formellement lisse (pour les topologies
discrètes), la suite
\begin{equation}
\tag{20.5.14.1}
\label{0.20.5.14.1-fr}
  0\longrightarrow\mathfrak{R}/\mathfrak{R}^2
  \xrightarrow{\delta_{C/B/A}}\Omega^1_{B/A}\otimes_B C
  \xrightarrow{u_{C/B/A}}\Omega^1_{C/A}\longrightarrow0
\end{equation}
est exacte et scindée.
\end{corollary}

\begin{proof}
En effet, toute $A$-extension de $C$ par un $C$-module est alors triviale
(19.4.4.1).
\end{proof}

